\begin{table}[!h]
\centering
\caption{Hospital Pre-Period Characteristics by Operating Margin "Win" Quartile}
\centering
\begin{tabular}[t]{lrrrr}
\toprule
Variable & Q1 (Smallest Win) & Q2 & Q3 & Q4 (Largest Win)\\
\midrule
N Hospitals & 339 & 339 & 338 & 338\\
Beds & 199 & 229 & 224 & 228\\
Operating Margin & 6.3\% & 0.5\% & -1.2\% & -0.1\%\\
Total Discharges & 10,208 & 11,975 & 11,671 & 11,965\\
Medicaid Share of Discharges & 15.7\% & 14.2\% & 14.1\% & 15.8\%\\
\addlinespace
Medicaid Discharges per Bed & 7.56 & 7.22 & 6.98 & 8.13\\
Uncompensated Care per Bed & \$125,424 & \$104,060 & \$83,608 & \$115,374\\
UCC-to-NPR Ratio & 14.8\% & 11.9\% & 10.4\% & 13.0\%\\
\% Rural & 9.5\% & 8.2\% & 12.5\% & 9.5\%\\
\% Nonprofit & 74.0\% & 87.3\% & 88.5\% & 78.4\%\\
\addlinespace
\% For-Profit & 26.0\% & 12.7\% & 11.5\% & 21.6\%\\
\bottomrule
\multicolumn{5}{l}{\textsuperscript{} Pre-period (rel\_year < 0) hospital characteristics, by quartile of the winsorized}\\
\multicolumn{5}{l}{individual hospital effect on operating margin (Q4 = largest improvement). See}\\
\multicolumn{5}{l}{heterogeneity\_win.R and heterogeneity\_analysis\_notes.md for the effect definition}\\
\multicolumn{5}{l}{and winsorizing rationale.}\\
\end{tabular}
\end{table}
